\section{Experimental Methodology}\label{sec:setup}
\subsection{Research Design}
The study uses a controlled, paired simulation experiment. Simulation was chosen over a physical testbed because the research questions require many repetitions of identical attacks against different defenses, with known attacker identities for scoring, which a physical deployment cannot provide at this scale. Within each comparison, a seed fixes the node layout, the attacker identities, and the mobility paths, and only the defense changes. One seeded 60-minute run is the experimental unit; nodes and windows within a run are dependent and are not counted as separate samples.

\subsection{Environment and Instruments}
We used Contiki-NG \cite{contiking2022} (commit \texttt{385d72a7658a}) with the Cooja simulator \cite{cooja2006} (submodule \texttt{fabff3ee}), native Cooja motes, RPL Lite with MRHOF \cite{rfc6719}, and CSMA with an always-on radio, inside a pinned Docker image. The only change to Contiki-NG is a patch that exposes message counters and three hooks (DIS input, DAO input, and forwarding). Analysis used Python~3.12 with NumPy, pandas, SciPy, and Matplotlib.

Each network has 30 nodes (a root and 29 clients) in a $200\times200$~m area with the root at the center. Placements are regenerated until the network is connected at 90\% of the 50~m radio range, so that no node is isolated by geometry. The unit-disk radio model has a 100~m interference range. In mobile runs, 30\% of the nodes follow the random waypoint model \cite{camp2002} at 0.5--2~m/s with 0--10~s pauses, a common setting for walking-speed industrial mobility. Clients send every 10~s with random jitter.

\subsection{Defenses Compared}
We compared seven configurations. Plain RPL without any defense shows the impact of each attack. Secure-RPL \cite{vermaranga2020} and RPL-MRC \cite{medjek2021} are the two DIS defenses that DIS-FDM was compared with, and DIS-FDM \cite{prajisha2025} itself is the base method. Li-MSD \cite{limsd2024} is the closest DAO defense, and we set its threshold to 94 DAOs per 1800~s, the largest count produced by an honest node in separate benign runs. The remaining two are AEF-RPL and AEF fixed, which is AEF-RPL with learning switched off. All baselines were implemented from their publications in the same firmware, so that they share the radio, routing, and traffic stack.

\subsection{Experimental Campaigns}
The main campaign (seeds 1--10, 1960 runs) contained the DIS-FDM reproduction (2, 5, and 10 attackers; DIS intervals of 3--30~s), DIS-FDM's window sensitivity, the calibration runs, and the main matrix. The main matrix crossed four attack modes (none, DIS, DAO, combined), three attacker counts (3, 7, and 14 of 29 clients, i.e.\ 10, 25, and 50\%), two mobility levels (0 and 30\%), and all defenses. A second campaign on fresh seeds (Table~\ref{tab:campaigns}) tested robustness and isolated each component. Using new seeds for these studies means that they replicate the main findings on networks that played no part in designing the method. Twenty further benign mobile runs (seeds 11--20, learning on and off) measured false isolation with no attacker present.

\begin{table}[t]
\caption{Second campaign on fresh seeds (740 runs).}\label{tab:campaigns}
\centering\footnotesize
\begin{tabularx}{\columnwidth}{@{}Xrr@{}}\toprule
Study & Runs & Seeds \\\midrule
Attack-rate robustness (DIS/DAO intervals 3/2, 10/10, 30/30~s) & 240 & 21--30\\
Matched attack-free references for ARI & 70 & per study\\
Component ablations (six variants) & 240 & 31--40\\
Reception loss (success 1.0, .95, .8) & 120 & 41--45\\
Callback CPU profiling & 60 & 41--45\\
Large-network feasibility pilots and timing traces & 10 & 1--2, 21\\\midrule
Total & 740 & \\\bottomrule
\end{tabularx}
\end{table}

The six ablation variants were full AEF-RPL, learning off, DAO fusion mass disabled, immediate DAO evidence (persistence of one window), route repair disabled, and calibration fixed from earlier benign seeds. Each removes one mechanism while keeping the rest of the pipeline working, which lets the study attribute an effect to a specific design choice.

\subsection{Variables and Measures}
Table~\ref{tab:variables} lists the variables. Detection is scored per run at the node level: a node counts as flagged if any legitimate observer isolates it. TPR is the fraction of attackers flagged, FPR the fraction of legitimate clients flagged, and latency the time from onset to the first flag of each attacker. Upward PDR is the fraction of client packets that reach the root, and downward PDR the fraction of received packets whose echo returns; the downward direction carries the DAO damage. ARI and DRT follow Section~\ref{sec:theory}. Control overhead is the total number of DIO messages, and radio power is estimated from transmit and listen times. Native callback CPU time measures the cost of the per-window procedure on the host.

\begin{table}[t]
\caption{Experimental variables.}\label{tab:variables}
\centering\footnotesize
\begin{tabularx}{\columnwidth}{@{}lX@{}}\toprule
Role & Variables \\\midrule
Independent & Defense configuration (evidence sources, decision layer, DAO handling)\\
Moderators & Attack mode, attacker fraction, mobility, attack interval, reception success\\
Dependent & TPR, FPR, precision, macro-F1, latency; upward and downward PDR, ARI, DRT; DIO count, radio power, CPU time\\
Controlled & Seed (layout, attackers, mobility), traffic, duration (3600~s), onset (300~s), radio ranges, MAC\\\bottomrule
\end{tabularx}
\end{table}

\subsection{Statistical Analysis}
Every cell has ten paired seeds in the main campaign. Because the per-seed differences are few and often skewed or tied, we used the Wilcoxon signed-rank test \cite{wilcoxon1945}, which needs no normality assumption, paired by seed. Following the proposal, tests were one-sided for H2 and two-sided for H3 and H4. $p$-values were adjusted with Holm's method \cite{holm1979} within each family, defined as one comparison and one metric across its six attacker-fraction and mobility conditions, and we report the matched-pairs rank-biserial correlation as the effect size. With ten pairs, the smallest attainable adjusted value in a six-test family is $p_H=.012$, reached when all ten differences have the same sign. The fresh-seed campaign reports paired mean differences with 95\% bootstrap intervals (10\,000 resamples) \cite{efron1994} and exact sign-flip tests, Holm-adjusted across every test in each study.

\subsection{Reliability and Validity}
To make the results reliable, every run stores its scenario, firmware hash, and compressed raw log. A replay script reparses all 2720 logs and rebuilds every table. The rebuilt main-campaign tables are byte-identical to the archived ones, and the second campaign shows zero differences. We also repeated three runs after moving the simulator to another container engine and obtained the original results.

For internal validity, pairing by seed removes topology and attacker placement as confounders. The detectors see only local observations. We audited the source code to confirm that the defense cannot read attacker identities or the attack schedule, and calibration finishes before any attack begins.

Construct validity concerns whether the metrics measure what we intend. Service is measured in three ways (conditional PDR, ARI, and DRT), since delivery ratios alone do not show nodes that have lost their route. False isolation is measured in attack-free runs, in which every flag is a false positive.

As for external validity, the simulated network uses an ideal radio and CSMA, and Section~\ref{sec:limitations} discusses what this means for physical deployments.
